\documentclass[handout]{beamer}
\mode<presentation>
\usepackage{xcolor}
\usepackage[toc]{multitoc}
\usepackage{color}
\usepackage{graphicx}
\graphicspath{{../buku_ajar/buku_ajar_logika_matematika/figures/}}
\usepackage{tikz}
\usetikzlibrary{shapes,arrows,positioning,calc}
\usepackage{listings}
\usepackage{lmodern}
\usepackage{amsmath}
\usetheme[secheader]{Boadilla}
\usecolortheme{whale}
\setbeamertemplate{caption}[numbered]
\setbeamertemplate{bibliography item}[text]
\setbeamertemplate{navigation symbols}{}

\newcommand{\logic}[1]{\textbf{\textit{#1}}}

\theoremstyle{definition}
\newtheorem{definisi}[theorem]{Definisi}
\theoremstyle{example}
\newtheorem{contoh}[theorem]{Contoh}

\renewcommand*{\multicolumntoc}{2}

\setbeamertemplate{footline}
{
  \leavevmode%
  \hbox{%
  \begin{beamercolorbox}[wd=.333333\paperwidth,ht=2.25ex,dp=1ex,center]{author in head/foot}%
    \usebeamerfont{author in head/foot}\insertshortauthor%
  \end{beamercolorbox}%
  \begin{beamercolorbox}[wd=.433333\paperwidth,ht=2.25ex,dp=1ex,center]{title in head/foot}%
    \usebeamerfont{title in head/foot}\insertshorttitle
  \end{beamercolorbox}%
  \begin{beamercolorbox}[wd=.233333\paperwidth,ht=2.25ex,dp=1ex,right]{date in head/foot}%
    \usebeamerfont{date in head/foot}\insertshortdate{}\hspace*{2em} \insertframenumber{} / \inserttotalframenumber\hspace*{2ex} 
  \end{beamercolorbox}}%
  \vskip0pt%
}

\subtitle[Logika Matematika]{Logika Matematika}
\title[Pemodelan Komputasi]{Pemodelan Masalah Komputasi Berbasis Logika Predikat}
\author{Cecep Suwanda, S.Si., M.Kom.}
\date{Maret 2025}
\institute{Universitas Bale Bandung}

\begin{document}

\maketitle

\section*{Daftar Isi}
\begin{frame}
\frametitle{Daftar Isi}
\tableofcontents
\end{frame}

% ============================================================
% SECTION 8-1: Basis Pengetahuan
% ============================================================
\section{Basis Pengetahuan}

\begin{frame}
\frametitle{Basis Pengetahuan (Knowledge Base)}
\begin{definisi}[Basis Pengetahuan]
\logic{Basis Pengetahuan} (KB) adalah himpunan kalimat logis (FOL) yang merepresentasikan pengetahuan tentang domain. Setiap kalimat berupa fakta atom (ground) atau aturan berkuantor (implikasi universal).
\end{definisi}
\end{frame}

\begin{frame}
\frametitle{Komponen KB}
\begin{enumerate}
  \item \textbf{Himpunan Fakta ($\Delta$)}: Pernyataan atomik spesifik.
  \item \textbf{Himpunan Aturan ($\Gamma$)}: Implikasi berkuantor universal.
\end{enumerate}
\begin{block}{Contoh}
$\Delta$: $Manusia(Marcus)$, $OrangTua(Hasan, Ahmad)$\\
$\Gamma$: $\forall x\, (Manusia(x) \rightarrow Mortal(x))$
\end{block}
\end{frame}

\begin{frame}
\frametitle{Komponen Pemodelan}
\begin{center}
\small
\begin{tabular}{|l|l|l|}
\hline
\textbf{Komponen} & \textbf{FOL} & \textbf{Contoh} \\ \hline
Entitas & Konstanta & $Marcus$, $Hasan$ \\ \hline
Properti & Predikat Arity-1 & $Manusia(x)$, $Laki(x)$ \\ \hline
Relasi & Predikat Arity $\geq 2$ & $OrangTua(x,y)$ \\ \hline
Aturan & $\forall$ implikasi & $\forall x\, (P(x) \rightarrow Q(x))$ \\ \hline
\end{tabular}
\end{center}
\end{frame}

\begin{frame}
\frametitle{Contoh Translasi}
``Marcus manusia'', ``Marcus penduduk Pompeii'', ``Semua penduduk Pompeii tewas 79''.\\
\textbf{Formal:}
\begin{itemize}
  \item $Manusia(Marcus)$, $Penduduk(Marcus, Pompeii)$
  \item $\forall x\, (Penduduk(x, Pompeii) \rightarrow Tewas(x, LetusanGunung, 79))$
\end{itemize}
\end{frame}

% ============================================================
% SECTION 8-2: Fakta, Predikat, Arity
% ============================================================
\section{Fakta dan Predikat}

\begin{frame}
\frametitle{Predikat Arity-1, 2, $n$}
\begin{block}{Arity-1 (Properti)}
$Manusia(x)$, $Laki(x)$, $Prima(n)$, $Login(u)$
\end{block}
\begin{block}{Arity-2 (Relasi Biner)}
$OrangTua(x,y)$, $Suka(x,y)$, $Teman(x,y)$. Urutan argumen penting!
\end{block}
\begin{block}{Arity-$n$}
$Tewas(x, penyebab, tahun)$; $Nilai(mhs, mk, semester, skor)$
\end{block}
\end{frame}

\begin{frame}
\frametitle{Fakta Ground vs Aturan}
\begin{center}
\small
\begin{tabular}{|l|l|l|}
\hline
\textbf{Aspek} & \textbf{Fakta} & \textbf{Aturan} \\ \hline
Variabel & Tidak ada & Variabel terikat \\ \hline
Kuantor & Tidak perlu & $\forall$ \\ \hline
Contoh & $Laki(Hasan)$ & $\forall x\,(Laki(x) \rightarrow \neg Perempuan(x))$ \\ \hline
\end{tabular}
\end{center}
\end{frame}

\begin{frame}
\frametitle{Contoh: Pemodelan Hobi}
\textbf{Fakta:} $Suka(Ali, Mat)$, $Suka(Budi, Info)$, $Suka(Ali, Info)$\\
\textbf{Aturan:} $\forall x,y,z\, (Suka(x,z) \wedge Suka(y,z) \wedge x \neq y \rightarrow Rekan(x,y))$\\
\textbf{Kesimpulan:} $Rekan(Ali, Budi)$ (sama-sama suka Informatika)
\end{frame}

% ============================================================
% SECTION 8-3: Aturan Inferensi Kuantor
% ============================================================
\section{Aturan Inferensi Kuantor}

\begin{frame}
\frametitle{Universal Instantiation (UI)}
$\forall x\, P(x) \Rightarrow P(c)$ untuk setiap $c \in \mathcal{U}$
\begin{contoh}
$\forall x\, (Manusia(x) \rightarrow Mortal(x))$, $Manusia(Socrates)$\\
UI: $Manusia(Socrates) \rightarrow Mortal(Socrates)$\\
Modus Ponens: $Mortal(Socrates)$
\end{contoh}
\end{frame}

\begin{frame}
\frametitle{UG, EI, EG}
\begin{block}{Universal Generalization (UG)}
$P(c)$ untuk $c$ sembarang $\Rightarrow$ $\forall x\, P(x)$
\end{block}
\begin{block}{Existential Instantiation (EI)}
$\exists x\, P(x) \Rightarrow P(c)$ untuk konstanta baru $c$ (witness)
\end{block}
\begin{block}{Existential Generalization (EG)}
$P(c) \Rightarrow \exists x\, P(x)$
\end{block}
\end{frame}

\begin{frame}
\frametitle{Ringkasan Empat Aturan}
\begin{center}
\small
\begin{tabular}{|l|l|l|}
\hline
\textbf{Aturan} & \textbf{Arah} & \textbf{Syarat} \\ \hline
UI & $\forall \Rightarrow P(c)$ & $c$ sembarang \\ \hline
UG & $P(c) \Rightarrow \forall$ & $c$ sembarang (tanpa asumsi) \\ \hline
EI & $\exists \Rightarrow P(c)$ & $c$ konstanta baru \\ \hline
EG & $P(c) \Rightarrow \exists$ & $c$ memenuhi $P$ \\ \hline
\end{tabular}
\end{center}
\end{frame}

% ============================================================
% SECTION 8-4: Resolusi Query
% ============================================================
\section{Resolusi Query}

\begin{frame}
\frametitle{Proses Resolusi Query}
\begin{enumerate}
  \item Terima query (mis. $Kakek(Hasan, Ali)$?)
  \item Cek apakah query ada di $\Delta$
  \item Cari aturan dengan konsekuen cocok
  \item Instansiasi aturan (UI)
  \item Evaluasi anteseden (rekursif jika perlu)
  \item Modus Ponens $\Rightarrow$ simpulkan
\end{enumerate}
\end{frame}

\begin{frame}
\frametitle{Forward vs Backward Chaining}
\begin{block}{Forward Chaining (Data-driven)}
Mulai dari fakta $\Delta$, terapkan aturan, turunkan fakta baru hingga query terjawab.
\end{block}
\begin{block}{Backward Chaining (Goal-driven)}
Mulai dari query, cari aturan dengan konsekuen cocok, buktikan anteseden secara rekursif.
\end{block}
\end{frame}

\begin{frame}
\frametitle{Contoh: Marcus Legenda}
$\Delta$: $Manusia(Marcus)$, $Penduduk(Marcus, Pompeii)$\\
$\Gamma$: $\forall x\,(Manusia \rightarrow Mortal)$, $\forall x\,(Penduduk(x,Pompeii) \rightarrow Tewas(...))$, $\forall x\,(Mortal \wedge Tewas \rightarrow Legenda)$\\
\textbf{Forward:} UI $\Rightarrow$ $Mortal(Marcus)$, $Tewas(...)$, $Legenda(Marcus)$ $\checkmark$
\end{frame}

% ============================================================
% SECTION 8-5: Silsilah Keluarga
% ============================================================
\section{Studi Kasus: Silsilah}

\begin{frame}
\frametitle{Pemodelan Silsilah}
\textbf{Fakta:} $OrangTua(x,y)$, $Laki(x)$, $Perempuan(x)$\\
\textbf{Aturan:}
\begin{itemize}
  \item $Ayah(x,y)$ jika $OrangTua(x,y) \wedge Laki(x)$
  \item $Kakek(x,z)$ jika $Ayah(x,y) \wedge OrangTua(y,z)$
  \item $Saudara(x,y)$ jika $\exists z\, (OrangTua(z,x) \wedge OrangTua(z,y) \wedge x \neq y)$
\end{itemize}
\end{frame}

\begin{frame}
\frametitle{Query: Kakek dari Ali?}
\textbf{Backward chaining} dari $R_3$: $Kakek(x,z)$ jika $Ayah(x,y) \wedge OrangTua(y,z)$.\\
$z=Ali$ $\Rightarrow$ $OrangTua(Ahmad, Ali)$ dari $\Delta$.\\
Cari $Ayah(x, Ahmad)$: $OrangTua(Hasan, Ahmad) \wedge Laki(Hasan)$ $\Rightarrow$ $Ayah(Hasan, Ahmad)$.\\
$\Rightarrow$ $Kakek(Hasan, Ali)$ \textbf{Benar}.
\end{frame}

% ============================================================
% SECTION 8-6: E-Commerce
% ============================================================
\section{Studi Kasus: E-Commerce}

\begin{frame}
\frametitle{Aturan Bisnis E-Commerce}
\begin{itemize}
  \item Login $\wedge$ Terverifikasi $\rightarrow$ BolehSimpan
  \item BolehSimpan $\wedge$ Simpan(item) $\wedge$ Saldo $\geq$ Harga $\rightarrow$ BolehCheckout
  \item BolehCheckout $\wedge$ Premium $\rightarrow$ Diskon 10\%
\end{itemize}
\end{frame}

\begin{frame}
\frametitle{Inferensi: UserA Diskon?}
$\Delta$: $Login(UserA)$, $Terverifikasi$, $Premium$, $Saldo(500k)$, $HargaItem(Laptop, 450k)$, $Simpan(Laptop)$\\
$R_1$ $\Rightarrow$ $BolehSimpan(UserA)$\\
$R_2$ (500k $\geq$ 450k) $\Rightarrow$ $BolehCheckout(UserA, Laptop)$\\
$R_3$ $\Rightarrow$ $Diskon(UserA, Laptop, 10\%)$ \textbf{Benar}
\end{frame}

% ============================================================
% SECTION 8-7: CWA, OWA, Batasan
% ============================================================
\section{CWA, OWA, dan Batasan}

\begin{frame}
\frametitle{CWA vs OWA}
\begin{definisi}[CWA]
Fakta yang \textbf{tidak} ada di KB dianggap \textbf{salah}. KB dianggap lengkap.
\end{definisi}
\begin{definisi}[OWA]
Ketiadaan fakta \textbf{tidak} berarti salah; status \textbf{tidak diketahui}. KB mungkin tidak lengkap.
\end{definisi}
\end{frame}

\begin{frame}
\frametitle{Perbandingan CWA dan OWA}
\begin{center}
\small
\begin{tabular}{|l|l|l|}
\hline
\textbf{Aspek} & \textbf{CWA} & \textbf{OWA} \\ \hline
Fakta tidak ada & Dianggap salah & Tidak diketahui \\ \hline
Contoh & SQL, basis data & Web Semantik, OWL \\ \hline
\end{tabular}
\end{center}
\begin{block}{Contoh}
Query $Laki(Budi)$? CWA: Salah. OWA: Tidak diketahui.
\end{block}
\end{frame}

\begin{frame}
\frametitle{Batasan FOL}
\begin{itemize}
  \item \textbf{Ketidakpastian}: Tidak menangani probabilitas (perlu fuzzy, Bayesian)
  \item \textbf{Perubahan}: Bersifat monoton; tidak ada \textit{defeasible reasoning}
  \item \textbf{Kompleksitas}: Semi-decidable; waktu komputasi signifikan
  \item \textbf{Temporal}: Tidak ada mekanisme bawaan untuk waktu
\end{itemize}
\end{frame}

% ============================================================
% RANGKUMAN
% ============================================================
\section{Rangkuman}

\begin{frame}
\frametitle{Rangkuman}
\begin{itemize}
  \item KB = Fakta ($\Delta$) + Aturan ($\Gamma$)
  \item Arity predikat: 1 = properti, 2+ = relasi
  \item UI, UG, EI, EG: instansiasi dan generalisasi kuantor
  \item Forward chaining (data-driven) vs Backward chaining (goal-driven)
  \item CWA: tidak ada = salah; OWA: tidak ada = tidak diketahui
\end{itemize}
\end{frame}

\begin{frame}
\frametitle{Kata Kunci}
\begin{center}
\small
basis pengetahuan $\bullet$ fakta $\bullet$ aturan $\bullet$ predikat $\bullet$ arity\\
UI $\bullet$ UG $\bullet$ EI $\bullet$ EG $\bullet$ forward chaining $\bullet$ backward chaining\\
CWA $\bullet$ OWA
\end{center}
\end{frame}

\end{document}
